\documentclass[10pt,a4paper]{amsart}
\usepackage[margin=19mm]{geometry}
\usepackage{amsmath,amssymb,mathtools}
\usepackage[scr=rsfs,cal=euler]{mathalfa}
\usepackage{xfrac}
\usepackage[unicode,hidelinks,bookmarks=false]{hyperref}
\newcommand{\ie}{i.e.\ }
\newcommand{\cf}{cf.\ }
\newcommand{\resp}{respectively\ }
\newcommand{\Subsection}{\subsection}
\providecommand{\nobreakdash}{\nobreak\hbox{-}}
\setlength{\emergencystretch}{2em}
\multlinegap0pt
\usepackage{bm}
\usepackage{bbm}
\usepackage{dsfont}
\usepackage{xspace}
\usepackage{cancel}
\usepackage{mathtools}
\usepackage{amsthm}
\makeatletter
\@ifundefined{theorem}{\newtheorem{theorem}{Theorem}}{}
\@ifundefined{proposition}{\newtheorem{proposition}[theorem]{Proposition}}{}
\makeatother
\newcommand{\N}{\mathcal N}
\newcommand{\cB}{\mathcal B}
\newcommand{\cG}{\mathcal G}
\newcommand{\cS}{\mathcal S}
\newcommand{\cT}{\mathcal T}
\title[Strong Subgraph-Count Stability in $C\_{2\ell+1}$-Free Graphs]{Strong Subgraph-Count Stability in $C\_{2\ell+1}$-Free Graphs}
\author{Yuanpei Wang, Xiamiao Zhao}
\date{Source: arXiv:2607.04347v1 (CC BY 4.0)}
\begin{document}
\maketitle

\noindent\emph{Source and licence.} Strong Subgraph-Count Stability in $C_{2\ell+1}$-Free Graphs. Version arXiv:2607.04347v1 (CC BY 4.0), distributed under the Creative Commons Attribution 4.0 International licence, as recorded in the arXiv licence metadata for this exact version and shown on the abstract page https://arxiv.org/abs/2607.04347v1. Full rights notice: licenses/arxiv-2607.04347-CC-BY-4.0.txt.\par\smallskip

\noindent\textit{Editorial scope.} Counted unit: the Proposition whose source label is \texttt{prop:path-extremal} in the article (source0.tex lines 318--324, source label \texttt{prop:path-extremal}), delivered together with the article's own complete original proof of it (source0.tex lines 326--371, one \texttt{proof} environment of 46 lines). No mathematics is written, repaired or re-derived: statement and proof are the article's own text line for line. Required context retained uncounted: none. Every definition, display and label of those lines is retained unchanged. Omitted from this package and not redistributed: 1--317, 325--325, 372--969. Nothing inside a retained excerpt is omitted or altered: every retained body is delivered exactly as the article's lines are written, so no omission receipt of any kind accompanies this package. Citations: the retained excerpts carry no citation command at all, so no bibliography entry of the article is transcribed and no reference is redistributed. Reference adaptation: none was needed and none was made; every reference inside the delivered unit resolves inside it, so no unresolved reference or citation is produced. Printed slips kept literal: no printed word, symbol, hypothesis or number of the counted statement or of its proof was altered. Layout and class: the article's own class is \texttt{article} with packages including fontenc, inputenc, lmodern, amsmath, amssymb, amsthm, mathtools, enumitem, microtype, hyperref, geometry; the wrapper is the lane's pinned \texttt{amsart} editorial wrapper at 10pt on A4, which restates the theorem-like environments the retained text needs and adds the article's own macro definitions that the retained text needs (\texttt{\textbackslash N}, \texttt{\textbackslash cB}, \texttt{\textbackslash cG}, \texttt{\textbackslash cS}, \texttt{\textbackslash cT}), with extra wrapper packages (bm, bbm, dsfont, xspace, cancel, mathtools). The delivered numbering is the unit's own. Assets: no retained span contains a figure environment, a graphics-inclusion command, a PDF-page inclusion, a table or a TikZ picture; no image file is copied into this package, the package declares no asset dependency and no image file is redistributed with it. Licence: version arXiv:2607.04347v1 of this article is distributed under the Creative Commons Attribution 4.0 International licence, as recorded in the arXiv licence metadata for this exact version and shown on the abstract page https://arxiv.org/abs/2607.04347v1. A full rights notice is delivered at licenses/arxiv-2607.04347-CC-BY-4.0.txt. Characters: every retained excerpt is pure ASCII, so the delivered engine, XeLaTeX with its default Latin Modern text font, reports no missing character.
\begin{proposition}\label{prop:path-extremal}
Fix integers $t\ge4$ and $r\ge3$. For all sufficiently large $n$,
\[
        \cB^{\mathrm P}_{t,r}(n)=\cB^{\mathrm P,\chi}_{t,r}(n)=\cS_{t,r}(n).
\]
Consequently, for every $S\in\cS_{t,r}(n)$, $b^{\mathrm P}_{t,r}(n)=b^{\mathrm P,\chi}_{t,r}(n)=\N(P_t,S)$.
\end{proposition}

\begin{proof}
Let $d=r-1$. The family $\cG_{n,r}$ is closed under deleting edges. Hence, if $G\in\cG_{n,r+1}\setminus\cG_{n,r}$ and $G^+$ is obtained from $G$ by adding edges with $G^+\in\cG_{n,r+1}$, then $G^+\notin\cG_{n,r}$. Also, in any expression of a graph in $\cG_{n,r+1}\setminus\cG_{n,r}$ by a bipartite core and suspended graphs, the number of outside vertices is $d$.

Let $G\in\cG_{n,r+1}\setminus\cG_{n,r}$ attain the maximum number of copies of $P_t$. Choose a bipartite core $B_0$ and suspended graphs for $G$ so that the outside set has size $d$, and let $A\cup B$ be a bipartition of $B_0$ with $|A|=a$ and $|B|=b$. Since $\cT^*(r,n)\subseteq\cG_{n,r+1}\setminus\cG_{n,r}$ and every member of $\cT^*(r,n)$ has $\Theta(n^t)$ copies of $P_t$, we must have $a,b=\Omega(n)$.

Let $\widehat G$ be obtained from $G$ by replacing $B_0$ with $K_{a,b}$ and replacing each suspended graph by the clique on its vertex set. Then $\widehat G\in\cG_{n,r+1}\setminus\cG_{n,r}$ and $G\subseteq\widehat G$. Since $a,b=\Omega(n)$ and $t$ is fixed, every edge in $E(\widehat G)\setminus E(G)$ is contained in a copy of $P_t$ in $\widehat G$. Indeed, an added cross-edge inside the core can be extended alternately in the two parts of $K_{a,b}$ until it has $t$ vertices. If an added edge joins the suspension vertex $x$ of some suspended graph to an outside vertex $u$, then one starts with $u x$ and continues from $x$ along an alternating path in the core; for $t=4$ this gives a path $u x y x'$ with $y$ in the opposite part of the core and $x'$ in the same part as $x$. If an added edge joins two outside vertices $u$ and $v$ in the same suspended graph, then one starts with $u v x$, where $x$ is the suspension vertex, and continues in the core; for $t=4$ this gives $u v x y$. Since the two core parts have linear size, the required distinct core vertices can always be chosen for all sufficiently large $n$. In each case the constructed copy contains the added edge and is therefore absent from $G$. Hence $G=\widehat G$ by the maximality of $G$.

We now maximize the number of copies of $P_t$ among the graphs obtained in this way. Such a graph is determined by $a+b=n-d$, by a partition $d=d_1+\cdots+d_s$ of the outside vertices into suspended cliques, and by the parts of the core containing the corresponding suspension vertices. For fixed data $\pi$, write $F_\pi(a,b)=\N(P_t,G_\pi(a,b))$ for the corresponding polynomial in $a$ and $b$.

Using the notation from the comparison above, let $d_i$ be the number of outside vertices in the $i$th suspended clique, and let $Z_i\in\{A,B\}$ be the side containing its suspension vertex. For $Z\in\{A,B\}$, set $Q_t^Z(a,b)=2R_{t-2}^Z(a,b)$. Since paths using at least three outside vertices, or two outside vertices from different suspended cliques, contribute only $O(n^{t-4})$ copies, we have
\[
\begin{aligned}
        F_\pi(a,b)
        &=\N(P_t,K_{a,b})
        +\sum_{Z_i=A} d_i R_{t-1}^A(a,b)
        +\sum_{Z_i=B} d_i R_{t-1}^B(a,b)\\
        &\quad
        +\sum_{Z_i=A}\binom{d_i}{2}Q_t^A(a,b)
        +\sum_{Z_i=B}\binom{d_i}{2}Q_t^B(a,b)
        +O(n^{t-4}).
\end{aligned}
\]
We first choose the sizes of the two parts of the complete bipartite core. Since $a+b=n-d$ is fixed, the number $\N(P_t,K_{a,b})$ is maximized when $a$ and $b$ are as equal as possible. More precisely, if $|a-b|\ge2$, then moving one vertex from the larger part to the smaller part increases $\N(P_t,K_{a,b})$ by $\Theta(n^{t-2})$. The terms involving outside vertices change by only $O(n^{t-3})$. Hence no extremal graph has $|a-b|\ge2$, and the core is $T_{n-d,2}=T_{n-r+1,2}$.

Now fix this balanced core. The side of the suspension vertex and the partition $d=d_1+\cdots+d_s$ must be compared together. Indeed, when the two parts of the core have sizes differing by one, the difference between $R_{t-1}^A(a,b)$ and $R_{t-1}^B(a,b)$ may have the same order as the terms involving two outside vertices.

For $Z\in\{A,B\}$ and $1\le j\le d$, define $M_Z(j)=jR_{t-1}^Z(a,b)+\binom j2 Q_t^Z(a,b)$.
Then the expansion above can be written as
\[
        F_\pi(a,b)=\N(P_t,K_{a,b})+\sum_{i=1}^s M_{Z_i}(d_i)+O(n^{t-4}).
\]
Let $\Lambda_Z=R_{t-1}^Z(a,b)+((d-1)/2)Q_t^Z(a,b)$, and choose $Z^*\in\{A,B\}$ so that $\Lambda_{Z^*}$ is maximal. Since $a,b=n/2+O(1)$, we have $Q_t^Z(a,b)=2R_{t-2}^Z(a,b)=c_{t,Z}n^{t-3}+O(n^{t-4})$ for some constant $c_{t,Z}>0$. Hence
\[
\begin{aligned}
        \sum_{i=1}^s M_{Z_i}(d_i)
        &=\sum_{i=1}^s d_i\left(R_{t-1}^{Z_i}(a,b)+\frac{d_i-1}{2}Q_t^{Z_i}(a,b)\right)\\
        &\le \sum_{i=1}^s d_i\left(R_{t-1}^{Z_i}(a,b)+\frac{d-1}{2}Q_t^{Z_i}(a,b)\right)\\
        &\le d\Lambda_{Z^*}=M_{Z^*}(d).
\end{aligned}
\]
If $s\ge2$, then each $d_i<d$, and the first inequality loses at least
\[
        \frac12\sum_{i=1}^s d_i(d-d_i)Q_t^{Z_i}(a,b)=\Omega(n^{t-3}),
\]
which dominates the error term $O(n^{t-4})$. Thus no extremal graph can have $s\ge2$. Hence all outside vertices lie in one suspended clique. The comparison for $\cT^*(r,n)$ above now gives precisely the members of $\cS_{t,r}(n)$. Each contains a clique $K_r$ and is $r$-colourable, so $\chi(S)=r$ for every $S\in\cS_{t,r}(n)$; hence the same graphs are extremal in the high-chromatic problem.
\end{proof}

\end{document}
